\documentclass[instructions.tex]{subfiles}
\begin{document}
 
\chapter{C’est ordinairement par notre faute que nous avons des ennemis et des chagrins.}
 
\primo Ayez soin de votre réputation, dit le Sage\footnote{Curam habe de bono nomine. Sap.}; c’est-à-dire, vivez de telle sorte que vous soyez sans reproche. Si vous vivez mal, c’est vous qui nuisez à votre réputation, et non pas les autres qui vous l’ôtent. Vous vous plaignez que vous avez des ennemis, vous devriez plutôt vous plaindre de vous-même. Si l’on ne vous aime pas, si l’on vous méprise, c’est que vous ne vous faites pas aimer. Les colombes ne font pas société avec le vautour, ni les oiseaux avec les hiboux, ni les agneaux avec les loups et les lions. Quelle société peut-on avoir avec vous, dont l’humeur inquiète et emportée, dont les airs fiers et impérieux, et dont les manières brusques et farouches rebutent tout le monde ?
 
Vous êtes un homme avide, intéressé, insensible aux misères d’autrui : quel amour, quelle estime peut-on avoir pour un homme dur et tenant, qui n’aime que soi-même, qui ne fait du bien à qui que ce soit ?
 
On parle mal de vous : pourquoi vous-même parlez-vous mal de tout le monde ? Avez-vous droit qu’on vous épargne, tandis que par vos satires et vos railleries vous n’épargnez personne ?
 
Vous vous plaignez qu’il n’y a plus de charité, que personne ne veut se fier à vous, ni vous rendre service : vous le méritez, parce que vous êtes un dissimulé, qui trompez tout le monde, qui n’opposez que la ruse et la chicane à la bonne foi, et l’ingratitude aux services qu’on vous rend. Plaignez-vous de vous seul, et non pas des autres. Vous êtes un ouvrier décrié : pourquoi vous décriez-vous vous-même par des ouvrages de mauvaise foi ? On vous abandonne dans votre pauvreté : c’est vous qui, par votre vie fainéante et par vos larcins, vous rendez indigne de la charité d’autrui.
 
On flétrit la réputation de vos enfans : pourquoi les exposez-vous à la censure des autres par les fréquentations que vous autorisez ? Et vous, filles du monde, pourquoi, par votre dissipation et vos airs mondains, vous attirez-vous la critique et les médisances du public ? Vivez dans la modestie, et vous ferez taire l’imposture.
 
\secundo Ceux qui nous haïssent et qui nous persécutent pèchent, parce qu’ils manquent de charité ; mais rendons-nous justice ; ne méritons-nous point, par notre conduite, leur aversion et leur censure ?
 
Que personne de vous, dit saint Pierre, ne souffre comme malfaiteur ; c’est-à-dire ne vous attirez aucune persécution par votre faute. Vous faites plus de tort aux autres, en leur donnant occasion de parler et de vous haïr, qu’ils ne vous en font par leurs persécutions; vous devez être plus affligé des péchés que vous faites commettre, en leur donnant lieu de manquer de charité, que des maux qu’ils vous font.
 
Si l’on vous persécute lorsque vous êtes innocent, réjouissez-vous d’être conforme à Jésus-Christ, qui a été le plus persécuté et le plus innocent des hommes.
 
Si l’on vous hait pour la justice, parce que vous faites votre devoir, que vous reprenez le vice, consolez-vous, c’est la destinée des bons pasteurs et des bons serviteurs de Dieu : Heureux, dit Jésus-Christ, ceux qui souffrent persécution pour la justice.
 
\section{Résolutions}
 
\primo Je m’efforcerai de ne m’attirer le mépris, la censure ni la haine de personne, tâchant de ne blesser aucunement les intérêts ni la délicatesse de qui que ce soit.
 
\secundo Ma conduite ne sera ni bizarre, ni incommode, ni vaniteuse ; et si je suis persécuté injustement, je m’en consolerai au pied de la croix. 

\prayer{O mon Dieu ! faites-moi la grâce de ne point m’attirer de chagrin par ma faute, et de supporter avec résignation et charité ceux qui m’arriveront d’ailleurs.}
 
\end{document}
